\section{A calibrated retiree}\label{sec:num}

The theorems predict exactly which of the two channels operates for which schedule, and a
calibration confirms it case by case. We solve a deterministic retiree of $25$ years by backward
induction with a global search over a grid of $6001$ asset levels --- the value function is not
concave, so a first-order condition cannot be used --- at $r=4\%$, $\beta=0.96$, relative risk
aversion $3$, a full pension of $0.40$ of mean earnings and a free area of $3$ mean-earnings units
(\texttt{numerics/meanstest.m}). The test is assessed on assets currently held, so all three
channels are live. The taper of $0.078$ per unit of assets is Australia's asset test, three dollars
a fortnight per thousand dollars.

\begingroup\small
\begin{center}
\begin{tabular}{@{}lccc@{}}
\toprule
Pension schedule & Cash on hand & Saving & Consumption \\
\midrule
Lump sum & increasing & increasing & increasing \\
Taper $0.078$ (Australia) & increasing & increasing & falls by $0.025$--$0.046$ \\
Taper $0.50$ & increasing & increasing & falls by $0.13$--$0.24$ \\
Taper $1.20$ (above $R$) & falls & falls by $0.004$ & falls by $0.55$--$0.70$ \\
Cliff & falls by $0.40$ & falls by $0.32$--$0.37$ & falls by $0.55$--$0.70$ \\
\bottomrule
\end{tabular}
\end{center}
\endgroup

Every entry is what the theorems require. Saving falls with assets in exactly the two rows where
cash on hand does, which are exactly the two rows whose taper exceeds the gross return, as
Theorem~\ref{thm:taper} says; in the other three it rises, as Theorem~\ref{thm:azmono} says it
must. Consumption falls in every row with a test of any kind, as Theorem~\ref{thm:cons} says it
can, and only in the lump-sum row is every policy monotone.

The consumption drop sits exactly at the free area and grows with age, from $0.025$ of mean
earnings at the start of retirement to $0.046$ twenty years in under the Australian taper. The
mechanism is the one the witness of Section~\ref{sec:cash} isolates: just above the free area the
effective return to saving is only $R-\tau$ until the pension is exhausted, so the household jumps
its saving towards the far side of the taper range and cuts consumption to pay for it. Under the
Australian schedule the borrowing constraint binds at about eight per cent of the state space at
every age.

The effect is not confined to retirees. A worker in the last year before retirement, facing the
Australian taper for the whole of retirement, has a monotone saving policy and a consumption policy
that falls by $0.025$ of mean earnings at the same threshold: the non-concavity is inherited
through the continuation value, and it is again consumption that carries it. That is
Theorem~\ref{thm:topkis} at work --- the continuation is as badly behaved as one likes, and saving
is monotone anyway.
